
\documentclass[12pt,a4paper]{article}
\usepackage[utf8]{inputenc}
\usepackage[T1]{fontenc}
\usepackage[spanish]{babel}
\usepackage{geometry}
\geometry{margin=2.2cm}
\usepackage{amsmath,amssymb}
\usepackage{booktabs}
\usepackage{float}
\usepackage{xcolor}
\usepackage{hyperref}

\title{\textbf{Informe Técnico de Analítica y Minería de Datos:\\ Reporte Ejecutivo Integral de Operación y Minería de Datos}}
\author{Sistema Inteligente de Monitoreo en Tiempo Real (Visión + Data Warehouse)}
\date{14/11/2026 19:06}

\begin{document}

\maketitle

\begin{abstract}
Este informe presenta los resultados obtenidos tras analizar 69 eventos en tiempo real mediante visión artificial, almacenamiento en Data Warehouse (DuckDB) y modelos de Minería de Datos. Se consolidó un total de 16 clientes atendidos y una ocupación media del 62.0\%.
\end{abstract}

\section{Enfoque Principal del Análisis}
Al no restringirse a una única variable, este informe sintetiza los hallazgos principales del sistema: se atendieron \textbf{16 clientes}, la mesa de mayor rotación fue \textbf{B1 (6p) (6 clientes, 2 ocupaciones)}, y el desempeño destacado del personal correspondió a \textbf{Mesero\_2 (1 atenciones, 5.18s promedio)}.

\section{Arquitectura del Sistema (Visión Artificial y Data Warehouse)}
El monitoreo se ejecutó mediante detección por visión artificial y rastreo determinista de entidades en salón, alimentando un \textbf{Data Warehouse dimensional en DuckDB} compuesto por la tabla de hechos \texttt{fact\_servicio\_live} (69 registros capturados) y las dimensiones \texttt{dim\_mesas} y \texttt{fact\_visitas\_meseros}.

\begin{table}[H]
\centering
\caption{Indicadores Globales Consolidados en el Data Warehouse}
\begin{tabular}{lr}
\toprule
\textbf{Indicador Operativo} & \textbf{Valor Registrado} \\
\midrule
Total de clientes atendidos hoy (con plato servido) & 16 \\
Clientes activos en el último corte & 13 \\
Mesas ocupadas en el último corte & 3 de 5 \\
Total de platos servidos detectados & 4 \\
Personal activo (Chefs / Meseros) & 4 Chefs / 2 Meseros \\
\bottomrule
\end{tabular}
\end{table}

\section{Desglose de Ocupación por Mesa (A1 -- B2)}
\begin{table}[H]
\centering
\caption{Rotación, Clientes Atendidos y Tiempo de Servicio por Mesa}
\begin{tabular}{lccccc}
\toprule
\textbf{Mesa} & \textbf{Sillas} & \textbf{Veces Ocupada} & \textbf{Clientes} & \textbf{Atenciones} & \textbf{Tiempo Prom.} \\
\midrule
A1 (4p) & 4 & 1 & 4 & 1 & 5.18 s \\
A2 (4p) & 4 & 2 & 0 & 0 & 0.00 s \\
A3 (4p) & 4 & 1 & 0 & 0 & 0.00 s \\
B1 (6p) & 6 & 2 & 6 & 0 & 0.00 s \\
B2 (6p) & 6 & 1 & 6 & 0 & 0.00 s \\
\bottomrule
\end{tabular}
\end{table}

\section{Productividad y Tiempos de Atención por Mesero}
\begin{table}[H]
\centering
\caption{Rendimiento y Tiempos Promedio del Personal de Salón}
\begin{tabular}{lcccc}
\toprule
\textbf{Mesero} & \textbf{Atenciones} & \textbf{Mesas Distintas} & \textbf{Tiempo Prom.} & \textbf{Tiempo Total} \\
\midrule
Mesero\_2 & 1 & 1 & 5.18 s & 5.2 s \\
\bottomrule
\end{tabular}
\end{table}

\section{Minería de Datos (Metodologías KDD, CRISP-DM y Modelo Híbrido)}
Siguiendo las fases de \textbf{KDD} (Selección, Preprocesamiento, Transformación, Minería y Evaluación) y \textbf{CRISP-DM}, se aplicó segmentación no supervisada mediante \textit{K-Means} y un \textbf{Modelo Híbrido Predictivo} (Árbol de Decisión + Red Neuronal Perceptrón Multicapa MLP) para anticipar saturación en el comedor:
\begin{itemize}
\item \textbf{Ocupación Media del Salón:} 62.0\% de capacidad instalada.
\item \textbf{Mesa Líder en Demanda:} B1 (6p) (6 clientes, 2 ocupaciones).
\item \textbf{Mesero Líder en Servicio:} Mesero\_2 (1 atenciones, 5.18s promedio).
\item \textbf{Modelo Híbrido (Árbol + MLP):} Evalúa la carga por mesero y variación de flujo para emitir alertas tempranas de saturación.
\end{itemize}

\section{Conclusiones y Recomendaciones Operativas}
El pipeline automatizado permitió transformar capturas visuales en inteligencia de negocio estructurada. Se recomienda priorizar la reposición rápida en \textbf{B1} y reforzar el salón cuando la ocupación supere el 60\% para mantener tiempos de atención óptimos.

\end{document}
